\begin{thebibliography}{99}
\footnotesize\itemsep=0pt

\bibitem{A-leyva}
Andruskiewitsch N., An introduction to {N}ichols algebras, in Quantization,
 Geometry and Noncommutative Structures in Mathematics and Physics, \textit{Math.
 Phys. Stud.}, \href{https://doi.org/10.1007/978-3-319-65427-0_4}{Springer}, Cham, 2017, 135--195.

\bibitem{A-nilpotent}
Andruskiewitsch N., On pointed {H}opf algebras over nilpotent groups,
 \textit{Israel~J. Math.}, {t}o appear, \href{https://arxiv.org/abs/2104.04789}{arXiv:2104.04789}.

\bibitem{AA}
Andruskiewitsch N., Angiono I., On {N}ichols algebras with generic braiding,
 in Modules and Comodules, \textit{Trends Math.}, \href{https://doi.org/10.1007/978-3-7643-8742-6_3}{Birkh\"auser}, Basel, 2008, 47--64,
 \href{https://arxiv.org/abs/math.QA/0703924}{arXiv:math.QA/0703924}.

\bibitem{AA-diag-survey}
Andruskiewitsch N., Angiono I., On finite dimensional {N}ichols algebras of
 diagonal type, \href{https://doi.org/10.1007/s13373-017-0113-x}{\textit{Bull. Math. Sci.}} \textbf{7} (2017), 353--573,
 \href{https://arxiv.org/abs/1707.08387}{arXiv:1707.08387}.

\bibitem{AAH-diag}
Andruskiewitsch N., Angiono I., Heckenberger I., On finite {GK}-dimensional
 {N}ichols algebras of diagonal type, in Tensor Categories and {H}opf
 Algebras, \textit{Contemp. Math.}, Vol.~728, \href{https://doi.org/10.1090/conm/728/14653}{Amer. Math. Soc.}, Providence,
 RI, 2019, 1--23, \href{https://arxiv.org/abs/1803.08804}{arXiv:1803.08804}.

\bibitem{AAH-triang}
Andruskiewitsch N., Angiono I., Heckenberger I., On finite {GK}-dimensional
 {N}ichols algebras over abelian groups, \href{https://doi.org/10.1090/memo/1329}{\textit{Mem. Amer. Math. Soc.}}
 \textbf{271} (2021), ix+125~pages, \href{https://arxiv.org/abs/1606.02521}{arXiv:1606.02521}.



\bibitem{AHS}
Andruskiewitsch N., Heckenberger I., Schneider H.J., The {N}ichols algebra of a
 semisimple {Y}etter--{D}rinfeld module, \textit{Amer.~J.~Math.} \textbf{132}
 (2010), 1493--1547, \href{https://arxiv.org/abs/0803.2430}{arXiv:0803.2430}.

\bibitem{ap1}
Andruskiewitsch N., Pe\~na Pollastri H.M., On the double of the (restricted)
 super {J}ordan plane, \textit{New York~J.~Math.} \textbf{28} (2022),
 1596--1622, \href{https://arxiv.org/abs/2008.01234}{arXiv:2008.01234}.

\bibitem{ASa-preNichols}
Andruskiewitsch N., Sanmarco G., Finite {GK}-dimensional pre-{N}ichols algebras
 of quantum linear spaces and of {C}artan type, \href{https://doi.org/10.1090/btran/66}{\textit{Trans. Amer. Math.
 Soc. Ser.~B}} \textbf{8} (2021), 296--329, \href{https://arxiv.org/abs/2002.11087}{arXiv:2002.11087}.

\bibitem{AS}
Andruskiewitsch N., Schneider H.J., Pointed {H}opf algebras, in New Directions
 in {H}opf Algebras, \textit{Math. Sci. Res. Inst. Publ.}, Vol.~43, \href{https://doi.org/10.2977/prims/1199403805}{Cambridge
 University Press}, Cambridge, 2002, 1--68, \href{https://arxiv.org/abs/math.QA/0110136}{arXiv:math.QA/0110136}.

\bibitem{angiono-campagnolo-sanmarco}
Angiono I., Campagnolo E., Sanmarco G., Finite {GK}-dimensional pre-{N}ichols
 algebras of super and standard type, \href{https://arxiv.org/abs/2009.04863}{arXiv:2009.04863}.

\bibitem{angiono-garcia}
Angiono I., Garc\'ia~Iglesias A., On finite {GK}-dimensional {N}ichols algebras
 of diagonal type: rank~3 and {C}artan type, \textit{Publ. Mat.}, {t}o appear, \href{https://arxiv.org/abs/2106.10143}{arXiv:2106.10143}.

\bibitem{angiono-garcia2}
Angiono I., Garc\'ia~Iglesias A., Finite {GK}-dimensional {N}ichols algebras of
 diagonal type and finite root systems, \href{https://arxiv.org/abs/2212.08169}{arXiv:2212.08169}.

\bibitem{brown-zhang}
Brown K.A., Zhang J.J., Survey on {H}opf algebras of {GK}-dimension~1 and~2, in
 Hopf Algebras, Tensor Categories and Related Topics, \textit{Contemp. Math.},
 Vol.~771, \href{https://doi.org/10.1090/conm/771/15507}{Amer. Math. Soc.}, Providence, RI, 2021, 43--62,
 \href{https://arxiv.org/abs/2003.14251}{arXiv:2003.14251}.

\bibitem{gr-jalg}
Gra\~na M., A freeness theorem for {N}ichols algebras, \href{https://doi.org/10.1006/jabr.2000.8363}{\textit{J.~Algebra}}
 \textbf{231} (2000), 235--257.

\bibitem{H-classif-RS}
Heckenberger I., Classification of arithmetic root systems, \href{https://doi.org/10.1016/j.aim.2008.08.005}{\textit{Adv. Math.}}
 \textbf{220} (2009), 59--124, \href{https://arxiv.org/abs/math.QA/0605795}{arXiv:math.QA/0605795}.

\bibitem{HS-adv}
Heckenberger I., Schneider H.J., Yetter--{D}rinfeld modules over bosonizations
 of dually paired {H}opf algebras, \href{https://doi.org/10.1016/j.aim.2013.05.009}{\textit{Adv. Math.}} \textbf{244} (2013),
 354--394, \href{https://arxiv.org/abs/1111.4673}{arXiv:1111.4673}.

\bibitem{HS-book}
Heckenberger I., Schneider H.J., Hopf algebras and root systems, \textit{Math.
 Surveys Monogr.}, Vol.~247, Amer. Math. Soc., Providence, RI, 2020.

\bibitem{KL}
Krause G.R., Lenagan T.H., Growth of algebras and {G}elfand--{K}irillov
 dimension, \textit{Grad. Stud. Math.}, Vol.~22, \href{https://doi.org/10.1090/gsm/022}{Amer. Math.
 Soc.}, Providence, RI, 2000.

\bibitem{R-quantum-groups}
Rosso M., Quantum groups and quantum shuffles, \href{https://doi.org/10.1007/s002220050249}{\textit{Invent. Math.}}
 \textbf{133} (1998), 399--416.

\bibitem{T}
Takeuchi M., Survey of braided {H}opf algebras, in New Trends in {H}opf Algebra
 Theory ({L}a {F}alda, 1999), \textit{Contemp. Math.}, Vol.~267, \href{https://doi.org/10.1090/conm/267/04277}{Amer. Math.
 Soc.}, Providence, RI, 2000, 301--323.

\bibitem{Ufer}
Ufer S., P{BW} bases for a class of braided {H}opf algebras,
 \href{https://doi.org/10.1016/j.jalgebra.2004.06.017}{\textit{J.~Algebra}} \textbf{280} (2004), 84--119, \href{https://arxiv.org/abs/math.QA/0311504}{arXiv:math.QA/0311504}.

\end{thebibliography}
